\chapter{Research Objectives}
\label{sec:research_objectives}

The primary aim of this PhD project is to improve informed decision-making in the diagnosis of ovarian cancer by updating and better understanding the \gls{ADNEX} model, ensuring that it can be used safely and effectively across diverse clinical settings. The project combines applied clinical research (Part~\ref{part:applied}) with methodological research (Part~\ref{part:methods}). Although the applied work centres on ovarian cancer diagnosis, the methodological insights are intended to be relevant more broadly to the development and evaluation of clinical prediction models.

This thesis is guided by five concrete research questions.

\begin{enumerate}
	\item What is the external validity of the \gls{ADNEX} model across all published validation studies, and which sources of bias or heterogeneity limit its use in practice? (Chapter~\ref{sec:chp_ADNEXSR})
	\item How does \gls{ADNEX} compare with the \gls{RMI} in direct head-to-head studies, both in discrimination and in clinical utility? (Chapter~\ref{sec:chp_ADNEXvsRMI})
	\item Can \gls{ADNEX} be updated using a more representative dataset, including patients who were not operated on, so that its predicted probabilities are better calibrated for clinical use? (Chapter~\ref{sec:chp_ADNEX2})
	\item How can model development and evaluation be improved so that predicted probabilities are more reliable, the uncertainty around individual risks is better understood, and overfitting is avoided? (Chapters~\ref{sec:chp_RFOverfitting} and~\ref{sec:chp_FundamentalProblem})
	\item How should calibration and clinical utility be evaluated when data are clustered, and can clinical utility be incorporated directly into variable selection? (Chapters~\ref{sec:chp_ClusteredCalibration} and~\ref{sec:chp_varselection})
\end{enumerate}

The expected outcome is a robust, well-calibrated, and clinically useful prediction model for ovarian cancer diagnosis with demonstrated utility, together with methodological advances that support informed decision-making across clinical prediction modelling more broadly.
